\section{Design}
\label{sec:design}

\begin{figure*}[t]
\centering
\resizebox{\linewidth}{!}{\begin{tikzpicture}[font=\small,
    pb/.style={draw,rounded corners=2pt,align=center,inner sep=5pt,minimum width=3.6cm,minimum height=1.2cm,thick},
    pl/.style={font=\footnotesize,align=center,fill=white,inner sep=1.5pt}]
  % columns at x = 0, 6.5, 13; rows at y = 0, -3.6, -7.0
  \node[pb,draw=untrusted,fill=untrusted!6] (vendor) at (0,0) {\textbf{Seller (x402)}\\writes price, payee,\\description};
  \node[pb,draw=gray,fill=gray!8] (user) at (7.2,0) {\textbf{User's console}\\checks the enrollment,\\opens and funds a goal};
  \node[pb,draw=enclave,fill=enclave!7] (signer) at (14.4,0) {\textbf{Signer}\\enrollment key,\\payer key per goal};
  \node[pb,draw=untrusted,fill=untrusted!6] (agent) at (0,-3.6) {\textbf{Orchestrator}\\LLM agent, relay key\\(gas only)};
  \node[pb,draw=chain,fill=chain!6] (vault) at (7.2,-3.6) {\textbf{EscrowPolicyVault}\\encrypted budget, cap;\\verifies attestations};
  \node[pb,draw=gray,fill=gray!8] (fac) at (14.4,-3.6) {\textbf{Facilitator}\\relays the signature\\to the vault};
  \node[pb,draw=gray!70!black,fill=gray!4,dashed] (inco) at (0,-7.0) {\textbf{Inco covalidators}\\off chain, in TEEs:\\evaluate and attest};
  \node[pb,draw=chain,fill=chain!6] (escrow) at (7.2,-7.0) {\textbf{Escrow, collector, USDC}\\accept calls only\\through the vault};

  % setup
  \draw[arr,enclave] (signer.west) -- node[pl,above]{\textbf{0a} enrollment} node[pl,below]{statement} (user.east);
  \draw[arr] (user.south) -- node[pl,pos=0.22,left=2pt]{\textbf{0b} open goal,\\fund payer} (vault.north);

  % one paid call
  \draw[arr] (vendor.south) -- node[pl,left=3pt]{\textbf{1} 402} (agent.north);
  \draw[arr] ($(agent.east)+(0,0.3)$) -- node[pl,above]{\textbf{2} request} ($(vault.west)+(0,0.3)$);
  \draw[arr] ($(agent.east)+(0,-0.3)$) -- node[pl,below]{\textbf{4} finalize} ($(vault.west)+(0,-0.3)$);
  \draw[darr] ($(vault.south west)+(0.4,0)$) -- node[pl,pos=0.45,above,sloped]{\textbf{3} events} ($(inco.north east)+(-0.4,0)$);
  \draw[arr] (inco.north) -- node[pl,left=2pt]{attested\\decision} (agent.south);
  \draw[arr,enclave] ($(agent.north east)+(-0.3,0)$) -- (1.5,-1.75) -- node[pl,pos=0.82,above]{\textbf{5} (goalId, seq)} (13.2,-1.75) -- ($(signer.south west)+(0.6,0)$);
  \draw[darr,enclave] ($(signer.south)+(0.6,0)$) .. controls (15.0,-1.5) and (10.5,-2.4) .. node[pl,pos=0.6,below]{reads record} ($(vault.north east)+(-0.4,0)$);
  \draw[arr] ($(signer.south)+(1.2,0)$) -- node[pl,right=2pt]{\textbf{6}\\signature} ($(fac.north)+(1.2,0)$);
  \draw[arr] (fac.west) -- node[pl,above]{\textbf{7} vault.authorize} (vault.east);
  \draw[arr] (vault.south) -- node[pl,right=2pt]{\textbf{8} fund (hold)} (escrow.north);
  \draw[arr] (vendor.east) -- (3.6,0) -- node[pl,pos=0.3,right=1pt]{\textbf{9} vault.claim} (3.6,-2.55) -- (5.9,-2.55) -- (5.9,-3.0);

  \begin{scope}[on background layer]
    \node[zone,draw=untrusted,fit=(vendor)(agent)] (zu) {};
    \node[zone,draw=chain,fit=(vault)(escrow)] (zc) {};
    \node[zone,draw=enclave,fit=(signer)] (ze) {};
  \end{scope}
  \node[font=\footnotesize\bfseries,text=untrusted,anchor=south] at (zu.north) {untrusted};
  \node[font=\footnotesize\bfseries,text=chain,anchor=north] at (zc.south) {on chain: decides, holds, pays};
  \node[font=\footnotesize\bfseries,text=enclave,anchor=south] at (ze.north) {signs};
\end{tikzpicture}
}
\caption{Architecture and one paid call. Red components read attacker-controlled text. Inco's covalidators
evaluate off chain and attest; the vault verifies the attestation on chain. The signer's only input from the
orchestrator is $(\mathit{goalId}, \mathit{seq})$, and its signature moves money only through the vault's
\code{authorize} (7, 8); the vendor is paid by \code{claim} through the vault (9).}
\label{fig:arch}
\end{figure*}

\subsection{Overview}

Figure~\ref{fig:arch} shows the components. The \emph{orchestrator} runs the language model and an x402
client, and is untrusted. The \code{EscrowPolicyVault} contract holds each user's policy, records every
request and decision, and is the operator of every escrowed payment. The \emph{signer} holds the enrollment
key and the payer keys. The user acts once per \emph{goal}: the console obtains an enrollment statement for a
fresh payer, checks it, sends \code{openEnrolledGoal} from the user's wallet with an encrypted budget, an
encrypted per-call cap, a public call count, a public attempt budget, an expiry and a payee allowlist, and,
once that transaction has a few confirmations, funds the payer with an amount chosen independently of the
budget (Section~\ref{sec:otherleak}). The user's wallet grants no allowance and never signs a payment.

A paid call proceeds as follows. The orchestrator receives a 402, parses it strictly, and lets the model decide
only whether to proceed; amount, payee and asset come from the parser. The relay key submits
\code{requestEscrowSpend(goalId, amount, payTo, resource, payBy, captureDeadline, keyHash)}. After Inco
evaluates the request, anyone may submit the attested decision to \code{finalizeDecision}. The orchestrator
asks the signer to authorize $(\mathit{goalId}, \mathit{seq})$ and attaches the signature to its retry. The
vendor's facilitator calls the vault's \code{authorize}, which funds the escrow, and the vendor then calls the
vault's \code{claim} to be paid. A paid call is four transactions, one of them the vendor's.
Figure~\ref{fig:seq}, after the references, gives every message.

\subsection{Authenticated enrollment}
\label{sec:enrollment}

The user funds the payer, and the payer's address reaches the console from the signer through the
orchestrator. Under A2 that path can substitute an address the attacker holds, and the user would fund it
directly: no request, policy or signature would be involved. Earlier versions of this system stated this as
a limit. Enrollment closes it.

\paragraph{Derivation.} The console draws a fresh 32-byte \emph{enrollment identifier} $e$ and asks for a
payer for its own wallet $o$. The signer derives the payer's key from the identifier
$\mathit{payer}/\mathit{chainId}/\mathit{vault}/o/e$ and returns the address with a statement
$\sigma_E$, an EIP-712 signature by its \emph{enrollment key} over $(\mathit{payer}, o, e)$, whose domain names
the chain and the vault. Because the payer is a function of $o$, no statement for a payer can ever name a
different owner: not one requested by a hostile orchestrator, and not one issued after a restart that lost the
signer's memory. The signer keeps no state about payers; the goal record names the owner and the vault records
$e$, so the key for any goal is re-derived from the chain. The enrollment key's identifier names only the
chain, so its address can be read before the vault that pins it is deployed.

\paragraph{Checks.} The vault fixes the enrollment address at deployment and accepts
\code{openEnrolledGoal}$(\mathit{params}, e, \sigma_E)$ only if $\sigma_E$ recovers to it for
$(\mathit{params.payer}, \mathit{msg.sender}, e)$; the unchecked \code{openGoal} reverts. The console, before
sending anything, checks the statement against the enrollment address built into it, compares that address
with the vault's \code{enrollmentSigner()} through its own RPC endpoint, and checks the owner, identifier, chain
and vault. Before funding it waits three blocks after the setup transaction and re-reads that the vault binds
this payer to this goal, that the goal is this wallet's, and that it was opened under $e$.

\paragraph{The connection to the measurement.} In a ROFL deployment the enrollment key is derived by Oasis's key
manager from the application's identity and the identifier \code{enrollment/<chainId>}, and released only to
instances the application's on-chain policy admits. That policy lists the accepted enclave measurements. A valid
statement therefore comes from an admitted instance, unless the key manager or the administrator who sets the policy
is compromised. On a laptop the key comes from a seed file, and a statement means only that whoever holds the file
signed it. We deployed the vault both ways and say for each run which one it used (Section~\ref{sec:liveeval}).

\paragraph{Finality of setup.} The adversary model allows dropped blocks, so a confirmed setup can disappear after
the user has funded the payer. The payer is then unbound, but every statement for it names this owner, and the vault
reads the owner from the sender, so no other wallet can bind it; the owner can open the goal again with the same
statement, and the funds wait in a payer only this owner's goal can use. Front-running fails for the same reason: a
statement seen in the mempool opens nothing for another sender. The confirmations narrow the window; they do not
carry the argument.

\subsection{Policy state and the debit}
\label{sec:writeahead}

For each goal the vault stores the remaining budget $B$ and the cap $C$ as encrypted handles, and the call count
$n$, attempt budget $m$, allowlist $A$, asset, expiry and relay in plaintext. The requested amount is public by
design: were it encrypted, the orchestrator could submit one ciphertext to the check and other terms to the signer.
Structural conditions (caller is the relay, goal open and unexpired, payee in $A$, non-zero amount, hold not passed,
an attempt left, no request pending) revert. Policy conditions (budget, cap and call count) resolve into a recorded
approval or rejection, so a refused request stays visible on chain.

The debit is written as a selection, since a contract cannot branch on an encrypted value:

\begin{lstlisting}[language=Solidity]
ok     = B[g].ge(a).and(C[g].ge(a));             // encrypted: a <= B and a <= C
debit  = ok.select(a.asEuint256(), uint256(0).asEuint256());
B[g]   = B[g].sub(debit);  B[g].allowThis();     // new handle, retained by the vault
ok.reveal();                                     // only the decision becomes public
\end{lstlisting}

All of this runs in the transaction that records the request, its sequence number $s$, a hash of its terms and a
frozen validity window. If the public call count is zero the vault records a rejection without touching the budget.
Both limits are combined in one condition, so a rejection does not say which limit refused. The call count moves at
finalization, and only on approval, which is why requests on a goal are strictly sequential.

This \emph{write-ahead} debit is a construction choice, not a security principle the bound needs. With one request
pending per goal, a design that debited at finalization, before any approval could be signed, would satisfy the same
invariants; ours is what Inco makes expressible, since the debit must be computed where the comparison is. What the
choice changes is the accounting: a pending request has already debited, so the invariants of
Section~\ref{sec:model} carry a term $Q$ for it. Selecting the operand, not the result, keeps the subtraction defined
on every path. Every debit produces a new budget handle to which only the vault is granted access.

\subsection{Finalization and the signer}
\label{sec:signer}

\code{finalizeDecision}$(g, s, v, \sigma)$ verifies that the covalidator signatures $\sigma$ attest plaintext $v$ for
\emph{the decision handle stored for $(g, s)$}, never one the caller supplies. Anyone may call it.

The signer accepts $(g, s)$ and nothing else. It reads the goal and the request, refuses unless the request is finalized,
approved, not reclaimed, on an open goal, in the expected asset and inside its frozen window, re-derives the payer's key
from the goal's owner and enrollment identifier, and refuses if that key is not the goal's payer. It then signs one kind of
message: a \code{receiveWithAuthorization} from the payer to the pinned collector, described in Section~\ref{sec:binding}.
The operator, collector, escrow and token come from configuration fixed at boot, never from the chain, and the signer has no
code path to a plain transfer. The record is frozen, so repeating the query gives a byte-identical signature. These refusals
tell a caller why nothing was signed; since the vault checks again when money moves, they are not what makes a signature safe.

After the owner closes a goal, the signer also signs a \emph{sweep} of the payer's remaining balance, payable to the goal's
owner, through the vault's \code{charge}, which checks on chain that the goal is closed, the receiver is its owner and no fee is
taken.

\subsection{Collection through the escrow}
\label{sec:escrow}

An EIP-3009 transfer signed for an approved request stays valid whatever later happens to the approval. In version~1 of this work
that put the signer's RPC inside the trusted base: an invented approval reported by a malicious endpoint would be signed, and a
signature released for an approval on a dropped branch survived the branch. With $B_0 = C = 0.20$ and $n_0 = 1$, an approval of
0.10 on a discarded branch and a canonical approval of 0.20 at the same sequence number carried different nonces, so the token
accepted both and paid 0.30. Waiting for finality was not practical: on 6 November 2026 Base Sepolia's \code{safe} block was about
23 minutes behind the head and its \code{finalized} block 3.8 hours, against a 60-second hold.

In the escrow design the vault is the operator of every payment. The vault's \code{authorize} derives the payment from its own record
and requires the caller's struct to equal it in every field, the record to be finalized, approved and not reclaimed, the goal to be open,
the amount to be the whole recorded amount, and the collector to be the pinned one; only then does it call the escrow. In the example
above, the 0.10 signature names a record that does not exist on the canonical chain, so it is refused and the payment is 0.20. The
inherited \code{requestSpend} reverts on this vault, so a goal has no other way to spend.

\begin{table*}[t]
\centering\small
\begin{tabular}{@{}L{0.3\linewidth}L{0.2\linewidth}L{0.44\linewidth}@{}}
\toprule
Escrow state & Time & \code{reclaimExpired} \\
\midrule
Not funded & before \textit{validBefore} & refused: the signature is still live \\
Not funded & from \textit{validBefore} & credit: token and escrow both refuse the signature from then on \\
Held & before capture deadline & refused: the vendor can still claim \\
Held & from capture deadline & void, funds to payer, credit \\
Funded, nothing held, nothing captured & any & credit: the payer took the funds back directly \\
Captured & any & refused: paid \\
\bottomrule
\end{tabular}
\caption{What decides a reclaim. The state is read from the escrow.}
\label{tab:reclaim}
\end{table*}

An approved request has three endings: paid, returned, or never funded. Table~\ref{tab:reclaim} gives the rule for crediting it back.
The vendor captures with \code{claim}$(g, s, k)$, which only the recorded payee may call and which requires a preimage of the request's
key hash if one was set. Closing a goal stops funding, because \code{authorize} requires an open goal; funds already held can be claimed
until their deadline and come back through a reclaim after it.

\subsection{The end-to-end binding}
\label{sec:binding}

Several contracts stand between a signature and a payment, and the loss bound depends on all of them. We state what binds what, for the
deployments our runs used.

\paragraph{The signed message.} The payer signs, under USDC's EIP-712 domain (name, version, chain id, token address),
$\code{ReceiveWithAuthorization}(\mathit{from} = \mathit{payer}, \mathit{to} = \mathit{collector}, \mathit{value} = \mathit{maxAmount},
\mathit{validAfter} = 0, \mathit{validBefore} = \mathit{preApprovalExpiry}, \mathit{nonce})$, where
\begin{align*}
h_0 &= \mathrm{keccak}(\mathrm{enc}(\mathit{TYPEHASH}, \mathit{info}[\mathit{payer} := 0])),\\
\mathit{nonce} &= \mathrm{keccak}(\mathrm{enc}(\mathit{chainId}, \mathit{escrow}, h_0))
\end{align*}
is the escrow's own payment hash with the payer zeroed, and $\mathrm{enc}$ is ABI encoding. The nonce therefore covers every field of
\code{PaymentInfo} except the payer (operator, receiver, token, maximum amount, the three deadlines, fee bounds, fee receiver and salt),
plus the chain and the escrow; the payer is covered by \textit{from}. The salt is $(g \ll 64) \mathbin{|} s$, which makes the nonce
unique per request and lets the vault find the record.

\paragraph{Who may submit it.} USDC accepts \code{receiveWithAuthorization} only from \textit{to}, the collector. The collector acts only
when called by the escrow, and checks that the nonce equals the escrow's payer-agnostic hash of the payment it is given. The escrow accepts
\code{authorize} and \code{charge} only from the payment's operator, the vault. So a signature can move tokens only inside the vault's
\code{authorize} or \code{charge}, and both re-derive the payment from the vault's record.

\paragraph{Every other route.} The escrow's operator-only functions are \code{authorize}, \code{charge}, \code{capture}, \code{void} and
\code{refund}. The vault calls \code{authorize} after its checks, \code{charge} only for a sweep to the owner of a closed goal, \code{capture}
only in \code{claim} for the recorded payee and the whole amount with no fee, and \code{void} only in a reclaim after the capture deadline. It
never calls \code{refund}. The payer's own \code{reclaim} returns a hold to the payer after the capture deadline, which Table~\ref{tab:reclaim}
counts. Because the vault never refunds and captures only the whole amount, the escrow's \code{refundableAmount} is non-zero exactly when a
payment was captured, which is what lets \code{reclaimExpired} tell ``paid'' from ``returned''. In another integration that called
\code{refund}, that inference would fail.

\paragraph{Versions and code identities.} The escrow is commerce-payments v1.1.0~\cite{commercepayments} (commit \code{6b0a5ef}), at
\code{0xf968\ldots{}B19c} on Base and Base Sepolia, with the ERC-3009 collector at \code{0x8612\ldots{}ab88}. Neither is a proxy. At Base Sepolia
block 56{,}480{,}000 the keccak hashes of their deployed code begin \code{0x33558f6c} and \code{0xe27322c8}, and the escrow's token-store
template's \code{0x8ff35d5d}. USDC is Circle's upgradeable proxy (code hash \code{0xedc5281a}), with implementation \code{0xd74c\ldots{}C5B5}
(\code{0x1254951e}) at that block; an upgrade by its administrator is outside our model. The repository's test fixture records the code and
every storage slot these contracts read, at that block.

\subsection{The recorded funding deadline}
\label{sec:hold}

x402 offers no reservation. We take the vendor's deadline $T$ to be the end of the payment window it declares: the orchestrator records when the
402 arrived, adds \code{maxTimeoutSeconds} and passes the result as \textit{payBy}. The vault sets
\[
  \mathit{validBefore} = \min(\mathit{payBy},\ t_{\mathrm{request}} + 1\,\mathrm{h},\ \mathit{expiry}),
\]
after which the token and the escrow both refuse to fund the payment, and records a capture deadline at or after \textit{validBefore}, at most one
hour later and not past the goal's expiry, after which the escrow refuses capture. What the implementation guarantees is therefore a
\emph{recorded funding deadline}: funding before \textit{validBefore}, capture before the capture deadline. It does not guarantee that
\textit{payBy} equals the vendor's own deadline (under A2 the orchestrator chooses it, within the hour), that capture happens before $T$, that the
vendor still holds what it offered, or that funding causes delivery. Even an honest orchestrator computes \textit{payBy} from a relative timeout, so
transit delay and a difference between the vendor's clock and the chain's need a margin.

Before committing, the orchestrator checks that enough of the hold remains for the request to be funded. The threshold is a heuristic, set from
measurement: in our runs the time from the request to funding had a maximum of \fundingmax\,s (Section~\ref{sec:eval}), and the check uses
\fundingcheck\,s. It is not a guarantee of success.

Approval does not move money. A crash, an abandoned request or a vendor that never claims can leave the budget reduced by money that was not paid.
\code{reclaimExpired}$(g, s)$ returns the amount to the budget and the call to the count in the cases of Table~\ref{tab:reclaim}. Anyone may call it, since
every condition is read from the chain. It needs finalization first, so a request whose decision Inco never reveals cannot be reclaimed; the owner's
close and sweep is then a separate recovery. Reclaim restores accounting and not necessarily the ability to spend: the goal must still be open, unexpired
and have an attempt left.

\subsection{The attempt budget}
\label{sec:attempts}

Every requested amount and every decision is public, so each request answers one question about the encrypted limits: is $a \le \min(B, C)$?
Whoever chooses the amounts chooses the questions. Rejections never consumed anything, and once reclaim existed an approval left to expire consumed
nothing either: name a \textit{payBy} one second away, never fund, reclaim. Our regression test recovers a secret cap of 137{,}391 units exactly in 23
requests by binary search, reclaiming every approval, so budget and call count end where they began. This is the general lesson of the paper's second
contribution: a mechanism that restores spending capacity also restores the ability to ask, and in a private policy asking is what leaks.

The goal therefore carries a public attempt budget $m \ge n_0$, set by the owner. Every request spends one attempt, approved or rejected, and nothing
returns it. The attempt is kept apart from the call: a call is a payment the owner agreed to fund and comes back if unused, an attempt is a question
already answered. The attempt budget is also a resource a vendor can exhaust: quotes above the limits make an honest agent spend attempts on
rejections. Section~\ref{sec:leakeval} measures the trade between leakage and honest completion.

\subsection{Execution traces}

Each run produces a hash-chained trace whose payment steps carry the decision handle, the covalidator signatures and the commit transaction;
rejected attempts stay in it. Only a Merkle root is stored on chain, once per (vault, goal). The anchor is permissionless and its first write is final,
so anyone can occupy a goal's slot with an arbitrary root; verification rejects that root, at a cost in audit availability only.
